\section{Software Modules}

\subsection{Capture Module}
The capture module is responsible for image acquisition on Raspberry Pi and represents the operational entry point of the whole data pipeline. Its implementation supports both manual one-shot execution and unattended periodic capture, allowing the system to operate in interactive and autonomous modes without architectural changes.

From a software perspective, this module encapsulates camera invocation parameters, runtime profile behavior, and output naming conventions. The naming strategy is not only cosmetic: it provides traceability and helps downstream modules distinguish capture origin and temporal order. The module is designed to work in coordination with service management logic, so scheduled or long-running capture can be activated without relying on manual shell interaction during routine operation.

An additional responsibility of this module is interoperability with optional external sources, such as ESP32-CAM streams accessed through backend proxy operations. Even in this mixed-source scenario, the capture module contributes to a unified lifecycle by maintaining coherent storage and retrieval behavior.

The same module also coordinates contextual acquisition by attaching best-effort GPS and BME280 readings to image metadata. This enables richer dataset semantics without making capture depend on mandatory sensor success.

\subsection{Gallery Module}
The gallery module provides the browsing and dataset management surface used by both web and tablet clients. Its core function is paginated retrieval of image entries enriched with operational metadata such as timestamp, file size, label status, annotation count, and contextual values (GPS/T/H/P). Pagination is implemented server-side and consumed uniformly by clients to maintain consistent behavior under large collections.

Beyond listing, the module handles item-level actions, including open, status refresh, single deletion, and bulk deletion workflows. It also integrates thumbnail-oriented rendering through lightweight server-side thumbnail endpoints to keep browsing responsive with large full-resolution datasets. These operations are implemented with explicit backend synchronization to avoid stale client-side assumptions. In practical terms, the module prioritizes coherence over optimistic local mutation, which is particularly important when multiple interfaces may act on the same dataset.

The gallery module therefore acts as both visualization layer and control layer for image lifecycle management, bridging raw acquisition output and annotation workflows.

\subsection{Labeling Module}
The labeling module is responsible for annotation authoring and persistence. It supports bounding-box creation, selection, editing, class assignment, and TP/FP state management. The interaction model is designed to remain effective in browser environments even when users need repeated micro-edits on dense images.

A key implementation property of this module is dual-format persistence. The JSON representation stores complete annotation state, enabling exact UI reconstruction and richer semantic status handling. In parallel, YOLO TXT export preserves compatibility with training pipelines and external tooling. The module ensures that these two artifacts remain logically aligned after every save operation.

Because annotation quality directly affects downstream model reliability, the module is designed to minimize ambiguity in save/load semantics and to expose clear feedback after updates. In this sense, it is both a UI module and a data integrity module.

\subsection{System Control Module}
The system control module exposes maintenance actions that are operationally critical in edge deployments. It includes server restart, full device reboot, and poweroff operations, all available through dedicated API endpoints and surfaced in client UI with explicit confirmation flows for high-impact actions.

This module reduces dependency on ad-hoc SSH intervention for routine maintenance, while preserving safety constraints. In particular, destructive actions are intentionally separated from frequent operational controls and require deliberate confirmation. The implementation thus balances accessibility with risk containment.

From an architectural standpoint, the module is also an example of edge-oriented design philosophy: device management is part of the application surface, not an external afterthought.

\subsection{Discovery Module}
The discovery module resolves reachable service endpoints for Raspberry Pi and, where applicable, ESP32-CAM devices. Its implementation follows a cache-first strategy: previously valid hosts are tested before launching wider subnet scans. This behavior significantly reduces perceived latency during repeated usage sessions.

Discovery is context-aware with respect to network mode. In AP scenarios, device probing behavior differs from standard Wi-Fi scenarios, and unnecessary scans are intentionally avoided when operational assumptions make them unlikely to produce useful results. This adaptive strategy improves responsiveness and reduces background network noise.

The module is tightly linked to usability: fast and reliable discovery determines whether control operations feel immediate or cumbersome. For this reason, discovery is treated as a dedicated module rather than embedded utility code.

\subsection{Module Interactions and Cohesion}
Although each module has a distinct responsibility, the implementation emphasizes cohesive interactions. Capture feeds gallery, gallery feeds labeling, labeling feeds export quality, and system/discovery modules preserve runtime operability. The backend API layer is the contract that coordinates these interactions and prevents coupling through hidden client-side assumptions.

This modular decomposition has two practical benefits. First, it allows targeted refactoring without destabilizing unrelated features. Second, it creates a clear path for future extensions, since new capabilities can be attached to existing contracts with controlled impact on already validated workflows.
